\section{Related Work}\label{sec:related}

\emph{Battle-royale prediction.} Early PUBG studies predict final placement from post-match aggregates such as the Kaggle placement data~\cite{kagglePubgFinishPlacement,rokad2019pubgPlacement,sen2021pubgFinalRank}, and Lee and Lee~\cite{lee2025pubgRankFactors} study factors that influence PUBG rank prediction; such inputs are complete only after the match. Rating-system evaluations~\cite{dehpanah2020freeForAllRating,dehpanah2021teamBattleRoyaleRating} likewise target ranking rather than in-match risk.

Sagong, Woo, and Kim~\cite{sagong2025hunterCamper} are the closest in-match work. Their unit is the player in solo matches only; duo and squad are excluded, citing the role of team strategy. At each phase the label is whether the player survives into the next half-phase (19 stages), with a winner label at the last stage. Eight per-phase accumulated features (FightScore and ActiveScore are unweighted sums) feed random forest, logistic regression, and XGBoost classifiers scored by accuracy and F1; the only signal about other players is a count of players within 100 m. Deep sequence models were set aside, the data are one week of confidential live-mode logs (7,719 matches), and hazard or Cox-type methods are not used. Their future work lists team-coordination variables for squad modes, sequence models, and other periods. Sagong and Kim~\cite{sagong2025survivalSequences} feed per-phase strategy sequences to recurrent models to predict the winner.
% VERIFY: only the abstract of Sagong and Kim (BigComp 2025, 2 pages) was
% checked; read the full text before submission, in particular whether a
% current-state baseline is reported and whether the unit is a player or a team.
Both motivate our phase axis.

KRAFTON AI released \texttt{pubg-winprob}~\cite{krafton_pubg_winprob}, a code repository for real-time win-probability prediction in PUBG esports matches. As implemented, a Transformer applies self-attention across up to 16 squads at sampled time points, counts enemies and teammates within 150 m, and has survival-time, Cox-hazard, and winner heads; a discrete-time hazard appears only as a baseline. The README refers to an accompanying paper, but we located no paper record as of October 2026, so we cite the software and describe only what the code shows. The code includes an attention-free encoder variant (\texttt{--encoder mlp}), so our RQ2 re-examines that mechanism under a different target (next-phase hazard), different data (public-API matches), and a different loss; we do not present it as a new test of cross-squad attention. Its data, splits, and checkpoints are not released, so our arms adapt its design rather than reproduce it.

\emph{Esports prediction and interaction graphs.} Win prediction has been studied for MOBA games~\cite{hodge2017mobaWinPrediction,yang2020honorKingsRealtime,junior2023lolRealtime}, and death prediction in Dota~2~\cite{katona2019dotaDeath} has a short-horizon event target similar to next-phase elimination. Graph models encode relations between teams or players explicitly~\cite{bisberg2022gcnWp,xenopoulos2021gnnSports}, and HIGFormer~\cite{wang2025higformer} builds a player--team heterogeneous interaction graph for soccer outcome prediction.

\emph{History versus current state.} The value of history depends on the setting. In Dota~2 death prediction, a 15-step GRU improved on a single-frame input, while 3-step and 30-step windows did not~\cite{ringer2023timeToDie2}. In ice hockey, an LSTM over traces of length 10 improved most metrics relative to length 1~\cite{liu2018hockeyIJCAI}. These mixed findings motivate an explicit current-state baseline (RQ1).

\emph{Discrete-time survival.} We cast elimination as a discrete-time hazard whose intervals are safe-zone phases~\cite{allison1982discrete,singer1993discrete}, evaluated with proper scores (log-loss and Brier); its likelihood over at-risk rows is a binary cross-entropy on interval-wise event indicators. Neural discrete-time models include a scalable hazard model~\cite{gensheimer2019nnetsurvival}, DeepHit~\cite{lee2018deephit}, and Dynamic-DeepHit~\cite{lee2020dynamicdeephit}, the last learning from longitudinal history. We use the plain hazard cross-entropy and do not propose a survival objective.

\emph{Time, agent, and relation-biased attention.} Factorised attention along time and agent axes is common in multi-agent trajectory modelling~\cite{ngiam2022scenetransformer,capellera2024transportmer}, and Graphormer~\cite{ying2021graphormer} adds graph structure as additive attention biases. We use these mechanisms as switches in one backbone to test information axes; the architecture itself is not the object of study. Within our search we did not find a battle-royale model that represents the shrinking safe zone as an explicit object node; RQ4 tests this representation against hand-crafted zone and habitat features.
